\documentclass[10pt,a4paper]{amsart}
\usepackage{amsmath,amssymb,amsthm,graphicx}
% Retained minimal preamble for the Stacks Project chapter more-algebra.tex (More on Algebra).
% The chapter file pinned-original-source.tex carries only \input{preamble} on its first line, so no
% preamble text lives inside the chapter itself.  The Stacks Project uses one single shared preamble for
% every chapter; the snapshot stages it identically as sources/stacks-chow/preamble.tex and
% sources/stacks-etale-cohomology/preamble.tex.  The block below is copied mechanically, line for line,
% from sources/stacks-chow/preamble.tex lines 14-18 (the Xy-pic setup) and 154-217 (the native theorem
% declarations together with the mathematical macros).  Nothing was added, removed or redefined; the
% wrapper class and the amsmath/amssymb/amsthm/graphicx base packages are supplied by the assembler.
% Chapter: more-algebra.tex -- More on Algebra.
% For commutative diagrams we use Xy-pic
\usepackage[all]{xy}

% We use 2cell for 2-commutative diagrams.
\xyoption{2cell}

% Theorem environments.
%
\theoremstyle{plain}
\newtheorem{theorem}[subsection]{Theorem}
\newtheorem{proposition}[subsection]{Proposition}
\newtheorem{lemma}[subsection]{Lemma}

\theoremstyle{definition}
\newtheorem{definition}[subsection]{Definition}
\newtheorem{example}[subsection]{Example}
\newtheorem{exercise}[subsection]{Exercise}
\newtheorem{situation}[subsection]{Situation}

\theoremstyle{remark}
\newtheorem{remark}[subsection]{Remark}
\newtheorem{remarks}[subsection]{Remarks}

\numberwithin{equation}{subsection}

% Macros
%
\def\lim{\mathop{\mathrm{lim}}\nolimits}
\def\colim{\mathop{\mathrm{colim}}\nolimits}
\def\Spec{\mathop{\mathrm{Spec}}}
\def\Hom{\mathop{\mathrm{Hom}}\nolimits}
\def\Ext{\mathop{\mathrm{Ext}}\nolimits}
\def\SheafHom{\mathop{\mathcal{H}\!\mathit{om}}\nolimits}
\def\SheafExt{\mathop{\mathcal{E}\!\mathit{xt}}\nolimits}
\def\Sch{\mathit{Sch}}
\def\Mor{\mathop{\mathrm{Mor}}\nolimits}
\def\Ob{\mathop{\mathrm{Ob}}\nolimits}
\def\Sh{\mathop{\mathit{Sh}}\nolimits}
\def\NL{\mathop{N\!L}\nolimits}
\def\CH{\mathop{\mathrm{CH}}\nolimits}
\def\proetale{{pro\text{-}\acute{e}tale}}
\def\etale{{\acute{e}tale}}
\def\QCoh{\mathit{QCoh}}
\def\Ker{\mathop{\mathrm{Ker}}}
\def\Im{\mathop{\mathrm{Im}}}
\def\Coker{\mathop{\mathrm{Coker}}}
\def\Coim{\mathop{\mathrm{Coim}}}

% Boxtimes
%
\DeclareMathSymbol{\boxtimes}{\mathbin}{AMSa}{"02}

%
% Macros for moduli stacks/spaces
%
\def\QCohstack{\mathcal{QC}\!\mathit{oh}}
\def\Cohstack{\mathcal{C}\!\mathit{oh}}
\def\Spacesstack{\mathcal{S}\!\mathit{paces}}
\def\Quotfunctor{\mathrm{Quot}}
\def\Hilbfunctor{\mathrm{Hilb}}
\def\Curvesstack{\mathcal{C}\!\mathit{urves}}
\def\Polarizedstack{\mathcal{P}\!\mathit{olarized}}
\def\Complexesstack{\mathcal{C}\!\mathit{omplexes}}
% \Pic is the operator that assigns to X its picard group, usage \Pic(X)
% \Picardstack_{X/B} denotes the Picard stack of X over B
% \Picardfunctor_{X/B} denotes the Picard functor of X over B
\def\Pic{\mathop{\mathrm{Pic}}\nolimits}
\def\Picardstack{\mathcal{P}\!\mathit{ic}}
\def\Picardfunctor{\mathrm{Pic}}
\def\Deformationcategory{\mathcal{D}\!\mathit{ef}}

% --- appended 2026-10-12 (Chief): the project own definitions of the reference/slogan/history
% display environments, copied mechanically from the shared Stacks preamble lines 8-12
% (\usepackage{verbatim} plus the three \newenvironment lines, which the chapter preamble keeps
% verbatim from the project preamble). Nothing added, removed or redefined: these environments
% carry no mathematical content and the project itself typesets them as comments. Without these
% three lines a counted statement whose own text uses them fails with Environment slogan undefined.
\usepackage{verbatim}
\newenvironment{reference}{\comment}{\endcomment}
\newenvironment{slogan}{\comment}{\endcomment}
\newenvironment{history}{\comment}{\endcomment}
\usepackage{cleveref}
\crefname{originallemma}{original lemma}{original lemmas}
\Crefname{originallemma}{Original lemma}{Original lemmas}
\makeatletter
\newlabel{lemma-pol-lifting}{{0BRF}{source}{}{sourcepointer-1}{}}
\newlabel{lemma-pol-lifting@cref}{{[originallemma][0][]0BRF}{[1][0][]source}{}{}{}}
\makeatother
\crefname{originallemma}{original lemma}{original lemmas}
\Crefname{originallemma}{Original lemma}{Original lemmas}
\makeatletter
\newlabel{lemma-pol-in-ideal}{{0H34}{source}{}{sourcepointer-2}{}}
\newlabel{lemma-pol-in-ideal@cref}{{[originallemma][0][]0H34}{[1][0][]source}{}{}{}}
\makeatother
\crefname{originallemma}{original lemma}{original lemmas}
\Crefname{originallemma}{Original lemma}{Original lemmas}
\makeatletter
\newlabel{algebra-lemma-universally-bijective}{{09EF}{source}{}{sourcepointer-3}{}}
\newlabel{algebra-lemma-universally-bijective@cref}{{[originallemma][0][]09EF}{[1][0][]source}{}{}{}}
\makeatother
\title{Polynomial identities and integrality for invariants of a finite group action modulo an ideal}
\author{The Stacks Project Authors}
\date{stacks-project:more-algebra.tex:local-snapshot-20261012; GNU FDL 1.2 or later}
\begin{document}
\maketitle
\setcounter{subsection}{1}
\noindent\textit{Editorial note.} All delivered mathematical text is the author's own, copied from the pinned Stacks Project chapter file more-algebra.tex; no statement or proof step was rewritten, reordered or removed. The counted span is the author environment and its own proof environment. References inside the span are handed to the assembler as navigation stubs that carry the canonical tag of the cited tag only; whether each of them is load-bearing is recorded on the stub itself.\par\medskip
\noindent\textit{Allowed standard prerequisites (editorial).}\begin{itemize}
\item Invariant elements b of (R/JR)\textasciicircum{}G lift to elements of R\textasciicircum{}G modulo the ideal generated by J, with a monic polynomial relation. invoked at proof line 32619. Supplies part (1).
\item For x in R\textasciicircum{}G the polynomial prod (T - sigma(x)) lies in the ideal generated by T - x and hence its reduction modulo J is determined by x mod J. invoked at proof lines 32620-32622. Gives (T - a)\textasciicircum{}n = T\textasciicircum{}n in part (2).
\item An integral ring map inducing purely inseparable residue field extensions induces homeomorphisms on spectra. invoked at proof lines 32624-32625. Gives the final spectral statement.
\end{itemize}\medskip
\ifdefined\hypertarget\hypertarget{sourcepointer-1}{}\fi
\noindent\textit{Source reference.} Original lemma 0BRF: sources/stacks-more-algebra/pinned-original-source.tex, source lines 32540--32540. Cited at proof line 32619 with I = JR for part (1); the conclusion is used inside the counted span.\par\medskip
\ifdefined\hypertarget\hypertarget{sourcepointer-2}{}\fi
\noindent\textit{Source reference.} Original lemma 0H34: sources/stacks-more-algebra/pinned-original-source.tex, source lines 32569--32569. Cited at proof lines 32622-32623 for the congruence (T - x)\textasciicircum{}n = T\textasciicircum{}n modulo J; the conclusion is used inside the counted span.\par\medskip
\ifdefined\hypertarget\hypertarget{sourcepointer-3}{}\fi
\noindent\textit{Source reference.} Original lemma 09EF: sources/stacks-chow/tags.txt, source lines 11898--11898. Cited at proof lines 32624-32625 for the spectral homeomorphisms; the conclusion is used inside the counted span.\par\medskip
\begin{lemma}
\label{lemma-invariants-modulo-bis}
Let $R$ be a ring. Let $G$ be a finite group of order $n$ acting on $R$.
Let $J \subset R^G$ be an ideal. Then $R^G/J \to (R/JR)^G$ is ring map
such that
\begin{enumerate}
\item for $b \in (R/JR)^G$ there is a monic polynomial
$P \in R^G/J[T]$ whose image in $(R/JR)^G[T]$ is $(T - b)^n$,
\item for $a \in \Ker(R^G/J \to (R/JR)^G)$
we have $(T - a)^n = T^n$ in $R^G/J[T]$.
\end{enumerate}
In particular, $R^G/J \to (R/JR)^G$ is an integral ring map which
induces homeomorphisms on spectra and purely inseparable extensions of
residue fields.
\end{lemma}

\begin{proof}
Part (1) follow from Lemma \ref{lemma-pol-lifting} with $I = JR$.
If $a$ is as in part (2), then $a$ is the image of $x \in R^G \cap JR$.
Hence $(T - x)^n = \prod_{\sigma \in G} (T - \sigma(x))$
is congruent to $T^n$ modulo $J$ by Lemma \ref{lemma-pol-in-ideal}.
This proves part (2). To see the final statement we may apply
Algebra, Lemma \ref{algebra-lemma-universally-bijective}.
\end{proof}

\end{document}
